% docs/documento_oficial/cap2_marco_teorico.tex
\section{Marco Teórico}

Este marco teórico presenta los fundamentos conceptuales y normativos que sustentan el desarrollo de DIAPCE. Se organiza en siete ejes: el diseño de mezclas de concreto y la ley de Abrams como base empírica, la influencia de las condiciones ambientales sobre el desempeño del concreto, los aditivos químicos y su clasificación, el desarrollo de resistencia a compresión a lo largo del tiempo, el modelado empírico adoptado por el sistema, los fundamentos tecnológicos de la solución y, finalmente, los trabajos relacionados que contextualizan la propuesta.

% ─────────────────────────────────────────────────────────────────────────────
\subsection{Diseño de mezclas de concreto estructural — Ley de Abrams y método ACI}

El diseño de mezclas de concreto es el proceso mediante el cual se determinan las proporciones de los materiales constituyentes ---cemento, agua, agregados finos y gruesos, y eventualmente aditivos--- con el fin de obtener un material que cumpla con los requisitos de resistencia, trabajabilidad y durabilidad especificados (ACI 211.1, 2016). El fundamento histórico de este proceso es la \textbf{ley de Abrams}, que establece una relación inversa entre la resistencia a compresión $f'_c$ y la relación agua-cemento (a/c). En forma log-lineal puede expresarse como:
\[
f'_c = K \cdot \left(\frac{w}{c}\right)^{-n}
\quad\text{o equivalentemente}\quad
\log f'_c = \log K - n \cdot \log\!\left(\frac{w}{c}\right),
\]
donde $K$ y $n$ son parámetros dependientes del tipo de cemento, los materiales y las condiciones de curado (Jiménez y Alvarado, 2020). A menor relación a/c, menor es la porosidad de la pasta cementante y mayor la resistencia resultante; sin embargo, la trabajabilidad de la mezcla disminuye, lo que puede comprometer la colocación y el acabado.

El método más difundido para traducir esta relación en proporciones prácticas es el del \textit{American Concrete Institute}, \textbf{ACI 211.1} (ACI 211.1, 2016), que define un procedimiento sistemático a partir de la resistencia de diseño ($f'_c$), el tamaño máximo del agregado grueso, el asentamiento requerido y la relación a/c. En Colombia, el diseño de mezclas estructurales también se rige por la \textbf{NSR-10} (Reglamento Colombiano de Construcción Sismo Resistente), cuyo título C establece los requisitos mínimos de resistencia y durabilidad y remite a los métodos del ACI. La verificación experimental de la resistencia se realiza conforme a la \textbf{NTC 673} y a la \textbf{ASTM C39} (ASTM C39, 2021), mediante ensayos de compresión sobre cilindros fabricados y curados según \textbf{ASTM C192} (ASTM C192, 2021).

En DIAPCE, la relación a/c opera como un nodo clave del flujo de decisión: el sistema restringe las opciones disponibles a los valores presentes en el conjunto experimental, reduciendo la incertidumbre del usuario al mostrar únicamente combinaciones observadas en laboratorio (Jiménez y Alvarado, 2020).

% ─────────────────────────────────────────────────────────────────────────────
\subsection{Condiciones ambientales y su efecto sobre el concreto}

El comportamiento del concreto durante su fraguado y endurecimiento está determinado en gran medida por las condiciones ambientales en las que se produce la reacción de hidratación del cemento. Las dos variables de mayor incidencia son la \textbf{temperatura} y la \textbf{humedad relativa} del entorno, razón por la cual DIAPCE las incorpora como los dos primeros filtros del flujo de decisión: temperatura $\rightarrow$ humedad $\rightarrow$ relación a/c $\rightarrow$ aditivo.

\subsubsection*{Temperatura}

La hidratación del cemento es una reacción exotérmica cuya velocidad aumenta con la temperatura. En climas cálidos ---como los que predominan en el departamento del Meta, donde las temperaturas medias oscilan entre 25 y 35~°C--- la hidratación se acelera, lo que puede resultar en una ganancia rápida de resistencia temprana pero en una resistencia final menor si no se controlan adecuadamente las condiciones de curado. Temperaturas superiores a 30~°C incrementan la tasa de evaporación del agua de mezcla, reducen el tiempo de trabajabilidad y pueden inducir fisuración plástica en las primeras horas. Por el contrario, en condiciones frías (por debajo de 10~°C), la hidratación se retarda significativamente, comprometiendo el desarrollo de resistencia a edades tempranas.

\subsubsection*{Humedad relativa}

La humedad relativa del ambiente afecta directamente la velocidad de evaporación del agua superficial del concreto. Una humedad relativa baja, combinada con viento o temperatura elevada, puede provocar una pérdida acelerada de agua antes de que el concreto alcance suficiente rigidez, generando fisuración por retracción plástica. El curado adecuado ---mediante láminas de polietileno, compuestos de curado o aspersión de agua--- busca contrarrestar esta pérdida y garantizar que la hidratación del cemento se complete en condiciones favorables (ASTM C192, 2021).

En el contexto del Meta, la combinación de temperatura elevada y humedad relativa alta en época lluviosa, o temperatura elevada con humedad baja en épocas secas, crea condiciones que se desvían de los supuestos implícitos en las tablas genéricas de diseño. Esto hace necesario fundamentar las decisiones de dosificación en datos experimentales producidos localmente, tal como propone DIAPCE.

% ─────────────────────────────────────────────────────────────────────────────
\subsection{Aditivos químicos para concreto}

Los aditivos químicos son materiales distintos del agua, los agregados y el cemento que se incorporan a la mezcla en pequeñas proporciones ---generalmente como porcentaje de la masa de cemento--- con el propósito de modificar alguna de sus propiedades en estado fresco o endurecido (ASTM C494, 2022). Su uso está regulado por la norma \textbf{ASTM C494/C494M} (ASTM C494, 2022) y por su equivalente colombiano \textbf{NTC 1299}. La selección guiada de aditivos en DIAPCE muestra únicamente los códigos válidos para la combinación temperatura–humedad–a/c seleccionada, garantizando coherencia con los datos observados.

\subsubsection*{Impermeabilizantes}

Los aditivos impermeabilizantes actúan reduciendo la permeabilidad del concreto endurecido, ya sea mediante la obturación de los poros capilares de la pasta cementante o mediante la formación de productos de reacción que cierran los canales de penetración de agua (ASTM C494, 2022). Su uso es especialmente relevante en estructuras expuestas a ambientes húmedos o suelos agresivos, condiciones frecuentes en el Meta dado su régimen de lluvias. En DIAPCE, los impermeabilizantes se representan mediante los códigos \textbf{P1}, \textbf{P2} y \textbf{P3}, correspondientes a dosificaciones del 2, 3 y 4~\% sobre la masa de cemento del producto \textit{Euco Vandex AM~10I} (marca Euco). El código \textbf{P0} representa la mezcla de control, sin aditivo.

\subsubsection*{Plastificantes}

Los aditivos plastificantes ---también denominados reductores de agua--- permiten mejorar la trabajabilidad de la mezcla sin incrementar la relación a/c, o bien mantener la trabajabilidad reduciendo el contenido de agua (ASTM C494, 2022). Son especialmente útiles en condiciones de temperatura elevada, donde la trabajabilidad se pierde rápidamente. En DIAPCE, los plastificantes se representan mediante los códigos \textbf{PP1}, \textbf{PP2} y \textbf{PP3}, correspondientes a dosificaciones del 0{,}2, 0{,}4 y 0{,}6~\% sobre la masa de cemento del producto \textit{Sika\textsuperscript{\textregistered} Plastiment\textsuperscript{\textregistered} AP} (marca Sika) (García y Torres, 2021).

La siguiente tabla resume la clasificación de aditivos adoptada en DIAPCE:

\begin{table}[H]
  \centering
  \caption{Clasificación de aditivos en DIAPCE (fuente: datos experimentales del semillero UCC)}
  \begin{tabular}{clll}
    \hline
    \textbf{Código} & \textbf{Tipo} & \textbf{Producto} & \textbf{Dosificación} \\
    \hline
    \textbf{P0}  & Control (sin aditivo) & N/A                       & 0\,\%   \\
    \textbf{P1}  & Impermeabilizante     & Euco Vandex AM 10I        & 2\,\%   \\
    \textbf{P2}  & Impermeabilizante     & Euco Vandex AM 10I        & 3\,\%   \\
    \textbf{P3}  & Impermeabilizante     & Euco Vandex AM 10I        & 4\,\%   \\
    \textbf{PP1} & Plastificante         & Sika\textsuperscript{\textregistered} Plastiment\textsuperscript{\textregistered} AP & 0{,}2\,\% \\
    \textbf{PP2} & Plastificante         & Sika\textsuperscript{\textregistered} Plastiment\textsuperscript{\textregistered} AP & 0{,}4\,\% \\
    \textbf{PP3} & Plastificante         & Sika\textsuperscript{\textregistered} Plastiment\textsuperscript{\textregistered} AP & 0{,}6\,\% \\
    \hline
  \end{tabular}
\end{table}

% ─────────────────────────────────────────────────────────────────────────────
\subsection{Desarrollo de resistencia a compresión}

La resistencia a compresión del concreto no se alcanza de forma instantánea: es el resultado de un proceso continuo de hidratación del cemento que se extiende durante semanas o incluso meses (ASTM C39, 2021). Por convención, la resistencia de diseño $f'_c$ se especifica y verifica a los \textbf{28 días}, edad a la que el concreto ha alcanzado aproximadamente el 100~\% de su resistencia de referencia bajo condiciones normales de curado (ACI 211.1, 2016). Las edades intermedias de \textbf{7 días} y \textbf{14 días} tienen valor diagnóstico: permiten anticipar si la mezcla sigue la curva esperada y tomar decisiones de curado o protección oportunas (ASTM C192, 2021).

La curva típica de desarrollo muestra un crecimiento rápido en los primeros días que se desacelera progresivamente: a los 7 días se alcanza aproximadamente el 65~\% de la resistencia de diseño; a los 14 días, alrededor del 85~\%; y a los 28 días, el 100~\% de referencia. Estas proporciones varían según el tipo de cemento, la relación a/c, la temperatura de curado y la presencia de aditivos (Jiménez y Alvarado, 2020). DIAPCE estima las resistencias $\{\hat{f}_7,\, \hat{f}_{14},\, \hat{f}_{28}\}$ como promedios empíricos del subconjunto de registros filtrado por condiciones ambientales, relación a/c y aditivo:
\[
\hat{f}_{t} \;=\; \operatorname{avg}\bigl\{\, f_{t,i} \;:\; (T_i,\, H_i,\, \mathrm{a/c}_i,\, \mathrm{aditivo}_i) \in \text{selección}\,\bigr\},
\quad t \in \{7,\, 14,\, 28\}.
\]
Este enfoque \textit{data-driven} privilegia la reproducibilidad y la coherencia con el comportamiento observado en el laboratorio regional.

% ─────────────────────────────────────────────────────────────────────────────
\subsection{Modelado empírico y enfoques de aprendizaje}

La literatura reciente explora desde regresiones multivariadas hasta redes neuronales para predecir $f'_c$ incorporando variables como a/c, tipo de cemento, granulometría, aditivos y condiciones de curado (Torres y Jiménez, 2023). DIAPCE adopta en esta versión un \textbf{modelo empírico de agregación} (promedios condicionados) por su:

\begin{itemize}
  \item \textbf{Transparencia}: resultados explicables y trazables a registros concretos del dataset (Jiménez y Alvarado, 2020).
  \item \textbf{Robustez} ante datos ruidosos cuando se realiza una estratificación adecuada por condiciones ambientales.
  \item \textbf{Compatibilidad} con las decisiones en cascada temperatura $\rightarrow$ humedad $\rightarrow$ a/c $\rightarrow$ aditivo y con las validaciones de rango (resistencia objetivo 24–57~MPa).
\end{itemize}

La arquitectura y el esquema de datos dejan abierta la integración futura de modelos de aprendizaje automático ---como regresión regularizada o \textit{tree ensembles}--- entrenados sobre las mismas vistas SQL, siguiendo la línea explorada por Torres y Jiménez (Torres y Jiménez, 2023).

% ─────────────────────────────────────────────────────────────────────────────
\subsection{Fundamentos tecnológicos de la solución}

\subsubsection*{Flutter y Dart como plataforma de desarrollo}

Flutter es un \textit{framework} de código abierto desarrollado por Google que permite construir aplicaciones nativas para Android, iOS, web y escritorio desde una única base de código escrita en el lenguaje Dart (Google LLC, 2024). Su arquitectura basada en \textit{widgets} y su motor de renderizado propio garantizan una experiencia visual consistente entre plataformas, lo que lo convierte en una opción adecuada para aplicaciones formativas que deben ejecutarse en los dispositivos heterogéneos de un entorno universitario (Tashildar et al., 2020). La interfaz de DIAPCE implementa formularios con validadores que restringen las entradas a los intervalos y opciones presentes en el dataset, y visualizaciones con \texttt{fl\_chart} para la composición de mezcla (gráfico de pastel) y la evolución de resistencia (gráfico de línea).

\subsubsection*{SQLite y almacenamiento local}

SQLite es un motor de base de datos relacional que opera directamente sobre el dispositivo, sin requerir un servidor externo (SQLite Consortium, 2024). Su ligereza y capacidad de funcionar sin conexión lo hacen ideal para entornos con conectividad limitada. En DIAPCE, SQLite gestiona el almacenamiento de parámetros de diseño, resultados de laboratorio, registros de usuarios y los datos experimentales del CSV, con normalización y llaves foráneas para garantizar integridad referencial, e índices para acelerar las consultas condicionales. El acceso se realiza mediante la biblioteca \texttt{sqflite} para Flutter (Tekartik, 2024).

\subsubsection*{Arquitectura limpia y datos experimentales en CSV}

El conjunto de datos que alimenta las recomendaciones de DIAPCE se almacena en un archivo CSV semicolon-delimitado (\texttt{csv\_datos\_1.1.csv}), generado a partir de los registros experimentales del semillero. Al iniciar la aplicación por primera vez, el archivo se procesa y sus registros se insertan en la base de datos SQLite local. La arquitectura del sistema sigue el enfoque de \textit{Clean Architecture}, separando las responsabilidades en capas: dominio (entidades y casos de uso), datos (fuentes locales, modelos) y presentación (interfaces de usuario).

% ─────────────────────────────────────────────────────────────────────────────
\subsection{Trabajos relacionados}

Diversas herramientas abordan el diseño de mezclas de concreto desde perspectivas complementarias. \textbf{ConcreBrain} utiliza redes neuronales para predecir la resistencia a partir de parámetros personalizados; \textbf{Concrete Mix Design} es una solución comercial basada en algoritmos de cálculo estándar; y \textbf{Concretelab} facilita el diseño y análisis de mezclas en entorno web. Frente a estas propuestas, DIAPCE se posiciona como un \textit{asistente data-driven} local con tres diferencias clave: opera completamente sin conexión, sus recomendaciones se derivan del conjunto experimental propio del semillero de la UCC ---no de modelos externos---, y su diseño es explícitamente formativo e institucional, no industrial. Como línea base normativa se emplean \textbf{ACI 211.1} (ACI 211.1, 2016), \textbf{ASTM C39} (ASTM C39, 2021), \textbf{ASTM C192} (ASTM C192, 2021) y \textbf{ASTM C494} (ASTM C494, 2022).
